\begin{landscape}
\subsection{Gaussian single-arm results}
\noindent Sc1 and Sc2 at $n_1=200$ include all five lrcBART priors and all four comparator settings.
\begingroup
\setlength{\tabcolsep}{4pt}
\renewcommand{\arraystretch}{1.13}
\renewcommand{\arraystretch}{1.13}
\fontsize{8.5}{10.3}\selectfont
\begin{longtable}{lrrrrrr}
\caption{Complete Gaussian single arm simulation results.}\label{tab:tableS6}\\
\toprule
Method / setting & Bias (MCSE) & RMSE (MCSE) & Post. SD & Width & Coverage (MCSE) & Power (MCSE) \\ \midrule
\endfirsthead
\multicolumn{7}{l}{\tablename\ \thetable\ (continued)}\\
\toprule
Method / setting & Bias (MCSE) & RMSE (MCSE) & Post. SD & Width & Coverage (MCSE) & Power (MCSE) \\ \midrule
\endhead
\midrule\multicolumn{7}{r}{Continued on next page}\\\endfoot
\bottomrule\endlastfoot
\multicolumn{7}{l}{\strut\textit{Mean difference: $n_1=200$, Sc1}}\\*
lrcBART: 100 & -0.015 (0.016) & 0.156 (0.012) & 0.176 & 0.689 & 0.950 (0.022) & 0.890 (0.031) \\
lrcBART: f90 & -0.012 (0.015) & 0.153 (0.011) & 0.163 & 0.636 & 0.960 (0.020) & 0.920 (0.027) \\
lrcBART: s0min & -0.015 (0.016) & 0.157 (0.012) & 0.151 & 0.591 & 0.910 (0.029) & 0.920 (0.027) \\
lrcBART: $w=0.9$, 100 & -0.010 (0.016) & 0.155 (0.011) & 1.289 & 6.036 & 1.000 (0.000) & 0.000 (0.000) \\
lrcBART: $w=0.9$, f90 & -0.007 (0.016) & 0.159 (0.010) & 1.292 & 6.067 & 1.000 (0.000) & 0.000 (0.000) \\
BART-CP & -0.014 (0.015) & 0.153 (0.011) & 0.150 & 0.588 & 0.940 (0.024) & 0.940 (0.024) \\
LM-CP & -0.029 (0.027) & 0.270 (0.017) & 0.252 & 0.983 & 0.960 (0.020) & 0.600 (0.049) \\
hier. LM & -0.025 (0.027) & 0.272 (0.017) & 0.576 & 2.210 & 1.000 (0.000) & 0.090 (0.029) \\
hier. LM & -0.026 (0.032) & 0.320 (0.025) & 1.621 & 6.617 & 1.000 (0.000) & 0.000 (0.000) \\
\multicolumn{7}{l}{\strut\textit{Mean difference: $n_1=200$, Sc2}}\\*
lrcBART: 100 & 0.062 (0.017) & 0.177 (0.010) & 0.179 & 0.697 & 0.980 (0.014) & 0.910 (0.029) \\
lrcBART: f90 & 0.064 (0.017) & 0.183 (0.010) & 0.184 & 0.714 & 0.970 (0.017) & 0.910 (0.029) \\
lrcBART: s0min & 0.061 (0.017) & 0.177 (0.010) & 0.168 & 0.653 & 0.980 (0.014) & 0.970 (0.017) \\
lrcBART: $w=0.9$, 100 & 0.065 (0.017) & 0.183 (0.010) & 1.262 & 5.920 & 1.000 (0.000) & 0.000 (0.000) \\
lrcBART: $w=0.9$, f90 & 0.065 (0.018) & 0.189 (0.010) & 1.268 & 5.950 & 1.000 (0.000) & 0.000 (0.000) \\
BART-CP & 0.056 (0.017) & 0.176 (0.010) & 0.168 & 0.657 & 0.930 (0.026) & 0.950 (0.022) \\
LM-CP & 0.290 (0.027) & 0.395 (0.024) & 0.273 & 1.067 & 0.850 (0.036) & 0.910 (0.029) \\
hier. LM & 0.357 (0.027) & 0.444 (0.024) & 0.588 & 2.243 & 1.000 (0.000) & 0.500 (0.050) \\
hier. LM & 0.293 (0.032) & 0.433 (0.030) & 1.605 & 6.502 & 1.000 (0.000) & 0.000 (0.000) \\
\end{longtable}
\noindent\begin{minipage}{\linewidth}\footnotesize Entries with parentheses are estimates (Monte Carlo standard errors). Post. SD is the mean posterior standard deviation, or the estimated standard error for PSCL; width is the mean 95\% interval width. Coverage and empirical power are proportions; Bayesian decisions require posterior probability of benefit above 0.95, whereas PSCL requires its lower 95\% confidence limit to exceed the benefit threshold. Each row has 100 replicates. Gaussian power uses a benefit threshold of 0.5; survival power uses a ratio threshold of 1. Targets 100, 75 and 50 are absolute counts; f90, f50 and f25 are fractions of the calibration ceiling. s0min fixes $s_0^2=10^{-6}$ without target selection. NP, CP and PP denote trial-only, completely pooled and source-adjusted control fits. In single-arm studies, CP uses all historical controls without a discrepancy term. The hierarchical variance prior is $\operatorname{IG}(3/2,3a/2)$. Within each scenario and sample-size block, the first and second hier. LM or hier. AFT rows use prior settings $a=0.05$ and $a=0.5$, respectively.\end{minipage}
\endgroup
\end{landscape}
